\documentclass[letterpaper,11pt]{article}
\usepackage[top=0.8in, bottom=0.8in, left=0.9in, right=0.9in]{geometry}
\usepackage[hidelinks]{hyperref}
\usepackage{lmodern}
\usepackage{parskip}
\usepackage{microtype}
\setlength{\parskip}{6pt}
\setlength{\parindent}{0pt}
\begin{document}
\begin{flushleft}
\textbf{\large Ajay Venkatesh} \\
Brooklyn, NY \\
+1 516 585 9013 \\
\href{mailto:av3855@nyu.edu}{av3855@nyu.edu}
\end{flushleft}
\vspace{10pt}
May 21, 2026
\vspace{6pt}

Hiring Manager \\
Progress Rail, A Caterpillar Company

\vspace{10pt}

Dear Hiring Manager,

My name is Ajay Venkatesh. I'm finishing my M.S. in Computer Engineering at NYU, with production software experience at PwC in Bengaluru, a summer SDE internship at Amazon, and ongoing on-campus work at NYU's Novel AI Technologies. I'm writing about the Business Data Analyst role at Progress Rail.

The role's emphasis on building data analytics tools that map relationships across Supply Chain, Operations, and Finance is where my engineering work has been concentrated. At PwC I spent the past few years on a distributed backend integrating with \textbf{Microsoft CRM} over \textbf{SQL Server}, writing complex queries, stored procedures, and parallelized batch ingestion that materially cut execution time. I built \textbf{Power BI} dashboards over datasets joined across multiple relational databases for business stakeholders, exactly the cost and performance analytics work the JD describes.

On data analytics development: at NYU's Novel AI Technologies I engineered \textbf{ETL pipelines} transforming \textbf{NoSQL} into \textbf{PostgreSQL} with concurrency control, then powered \textbf{AWS QuickSight} dashboards for real-time analytics. My Big Data graduate project built a \textbf{PySpark / HDFS} pipeline over millions of flight records, training Random Forest and Gradient Boosted Tree models with EDA, anomaly detection, and feature engineering, useful patterns for supply-chain analytics where irregular data and operational metrics need to be turned into actionable insights.

On stakeholder coordination: I've owned features end-to-end alongside architects, product managers, and business stakeholders at PwC, translating ambiguous operational questions into reporting and analytical deliverables for non-technical audiences across Operations and Finance teams.

Stack-wise: \textbf{SQL}, \textbf{Python}, \textbf{Power BI}, \textbf{Tableau}, plus \textbf{PostgreSQL}, \textbf{SQL Server}, \textbf{MySQL}, \textbf{Snowflake}, and \textbf{QuickSight} for visualization. Comfortable with the data-prep-to-dashboard pipeline that operations analytics depends on.

What pulls me toward Progress Rail is the operational seriousness. Supply chain and operations analytics for railroad equipment runs on accuracy that translates into real cost and performance outcomes, where careful, documented work matters more than dashboard novelty.

I'd welcome the chance to discuss further. Thanks for considering my application.

\vspace{12pt}
Sincerely, \\
Ajay Venkatesh

\end{document}